% Standalone research note. Intentionally supplied as uncompiled TeX.
\documentclass[11pt]{article}
\usepackage[T1]{fontenc}
\usepackage[utf8]{inputenc}
\usepackage[margin=1in]{geometry}
\usepackage{amsmath,amssymb,amsthm,mathtools,enumitem}
\usepackage[hidelinks]{hyperref}
\newtheorem{theorem}{Theorem}[section]
\newtheorem{proposition}[theorem]{Proposition}
\newtheorem{lemma}[theorem]{Lemma}
\newtheorem{corollary}[theorem]{Corollary}
\theoremstyle{definition}
\newtheorem{definition}[theorem]{Definition}
\newtheorem{example}[theorem]{Example}
\theoremstyle{remark}
\newtheorem{remark}[theorem]{Remark}
\newcommand{\Set}{\mathbf{Set}}
\newcommand{\Cat}{\mathbf{Cat}}
\newcommand{\Preord}{\mathbf{Preord}}
\newcommand{\FinMon}{\mathbf{FinMon}}
\newcommand{\FinGrp}{\mathbf{FinGrp}}
\newcommand{\FinSet}{\mathbf{FinSet}}
\newcommand{\V}{\mathcal V}
\newcommand{\W}{\mathcal W}
\newcommand{\E}{\mathcal E}
\newcommand{\C}{\mathcal C}
\newcommand{\N}{\mathbb N}
\newcommand{\Op}{\operatorname{Op}}
\newcommand{\Nat}{\operatorname{Nat}}
\newcommand{\Core}{\operatorname{Core}}
\newcommand{\Str}{\operatorname{Str}}
\newcommand{\ev}{\operatorname{ev}}
\newcommand{\im}{\operatorname{im}}
\newcommand{\id}{\operatorname{id}}
\newcommand{\cost}{\operatorname{cost}}
\newcommand{\supp}{\operatorname{supp}}
\newcommand{\eps}{\varepsilon}
\title{Native Operations Across Semantic Doctrines\\
\large Recognition, finite extraction, and the limits of semantic reconstruction}
\author{}\date{5 November 2026}
\begin{document}
\maketitle
\begin{abstract}
A category of models and its observable functor reconstruct natural semantic
operations. It does not, without further hypotheses, identify those operations
with programs in a chosen native language. Existing enriched
structure--semantics theorems already solve the recognition problem when
suitable relative free arities are admitted and their arrows match that
language. This note makes the residual problem precise when the tests omit
those free objects or the semantic doctrine changes.

The main constructive result concerns finite algebraic grammars. A finite
arity probe with a singleton fibre recognizes a native term among all natural
operations. Under a primitive interpretation, a target term has a source
implementation exactly when its value belongs to a generated subalgebra of
one such finite probe. If reducts of target tests cover the source tests under
finite products, subalgebras and quotients, native realizability is preserved
and reflected, with explicit source-term extraction and coherent composition
of comparisons. These are independently stated probe and model hypotheses,
not an assumption that every semantic operation already lifts.

Finite monoids give a complete positive calculation and a sharp limit:
semantic operations are profinite words, while finite words are recognized by
truncation probes. The idempotent power is not a monoid term but is uniformly
computable with finite-table access and iteration. A genuine change from
category-valued witnesses to reachability can additionally create operations
with no natural source lift. Finite native-table constraints and an exact
free-path section criterion give positive witness-lifting classes. The
project's existing law-and-possibility calculus supplies a related, already
proved change-of-doctrine result with generally noninvertible comparisons.

The result is a coherent framework with proved sufficient criteria and
explicit boundaries, not a universal reconstruction of executable structure
or a claim of a new general discovery. Published generic-model,
profinite, invariant-relation and change-of-base mechanisms are credited,
and empirical access, source-program generation, resources and native
witnesses remain distinct obligations.
\end{abstract}
\tableofcontents

\section{The question and the level of the answer}

The desired reconstruction is stronger than recovering a quotient of a fixed
family of linear operators. One would like to start from an independently
specified operational presentation, reconstruct the operations justified by
its semantics, and compare that reconstruction when the admissible models,
arities, evidence or execution doctrine changes. The central diagram is
\[
  \mathsf P_d\longrightarrow\mathsf N_d
       \xrightarrow{e_d}\Op_d.
\]
Here $\mathsf P_d$ is a supplied source grammar, $\mathsf N_d$ is an
ambient category of native realizations and their actual witnesses, and
$\Op_d$ is reconstructed from semantic presentations. The first arrow
need not generate all of $\mathsf N_d$. The second may identify programs,
erase execution histories, or forget possible outcomes. Even when their
composite is onto operations, it need not be onto witnesses or preserve a
specified resource bound. These levels will not be silently identified.

Four questions consequently have different answers.
\begin{enumerate}
\item \emph{Semantic reconstruction.} Which operation families commute
with all admitted transformations of models?
\item \emph{Native recognition.} Which such families are represented by
one expression in the declared source grammar?
\item \emph{Transport.} Does the actual comparison between doctrines
preserve or reflect that property, including the intended resources?
\item \emph{Witness recovery.} Can realizations or certificates be chosen
naturally and compatibly with composition, rather than separately at each
model or each arrow?
\end{enumerate}
A fifth, empirical question asks whether the native primitives and probes
can actually be prepared and executed on a physical source. None of the
first four alone answers it.

The grammar must specify sorts, constructors, equations, structural rules,
resources and access. A multiplication expression, a program with equality
and unbounded iteration, and a finite lookup table are different grammars.
An ambient morphism is not automatically a source program. Throughout,
``native'' means native relative to the grammar explicitly named in the
statement. When only ambient realizability is proved, that is said instead.

\subsection{What already works and what is added here}
The earlier CALF strand reconstructs residual states and a syntactic monoid
from a complete intervention response; bounded exact identification requires
its separate state and access promises. The subsequent
\emph{Operational Theory Formation Through Interaction} already gives
native-image recognition, finite bases, normal forms, executable refutations
and theory/monad change for its certified finite linear regime. It is a
positive theorem, not merely a warning about semantic maps.

The present note changes the question from identifying the quotient of
already supplied operators to recognizing a generated native repertoire
inside operation families natural across a category of presentations. Its
main additional proved statement is source-term extraction under a primitive
grammar change from an independently isolating finite probe, with finite
model coverage making that extraction conservative. The categorical
comparison is also allowed to change the \emph{kind of evidence}, and not
only the equations of one fixed linear base. The examples are organized as
tests of this one native-to-semantic comparison: completion can add
operations; a doctrine change can collapse native distinctions or create
unliftable operations; a witness lift needs its own coherence criterion.

This organization does not establish historical novelty. Free-object
recognition, enriched structure--semantics, implicit operations, isolated
finite words and native lifting in the project's O7 calculus are existing
positive results. The proofs below make particular finite-recognition and
transport consequences explicit. Whether those consequences, in this exact
packaging, constitute a new theorem in the literature is left open.

\section{Semantic reconstruction and correctly typed doctrine comparison}
\label{sec:framework}
A semantic doctrine supplies a symmetric monoidal closed category $\V$,
$\V$-categories $\E,\C$, a $\V$-functor $U:\E\to\C$, and a supplied
family of arities $j:\mathcal J\to\C$. The objects of $\E$ are the
admitted presentations or environments and its morphisms are the admitted
presentation transformations. Their choice is part of the evidence policy,
not something recovered from objectwise observations. Write
\[
 H_a=\C(ja,U-):\E\longrightarrow\V.
\]
Assume the following particular enriched hom-objects exist and are small
in the chosen universe:
\begin{equation}\label{cat-eq:op}
 \Op_d(a,b)=[\E,\V](H_a,H_b).
\end{equation}
We use the input-to-output convention: an element of this hom-object is
an operation $H_a\Rightarrow H_b$. Composition of transformations gives
a $\V$-category on the arity labels. The variance of $H$ makes
$\mathcal J^{\mathrm{op}}$ map to it. Merely forming these hom-objects
need not produce a $\mathcal J$-theory in the stronger technical sense.
For that identification one must impose the requisite tractability and
representability conditions. In particular, Proposition 4.24 of
\cite{LWP} gives the tractability criterion; its structure--semantics
adjunction is Theorem 4.27, not Proposition 4.24.

A native repertoire is supplied separately. At minimum it consists of a
typed grammar $P_d(a,b)$, an interpretation
\[
 e_d:P_d(a,b)\longrightarrow\Op_d(a,b)_0,
\]
and resource classes $P_d^{\le r}(a,b)$ whose union is the grammar.
Here the subscript $0$ means global elements, not an assertion that an
enriched hom-object is discrete. Typing and native substitution must be
interpreted compositionally. A category-valued repertoire additionally
supplies a grammar of native proof or execution cells with endpoints and
composition. Neither the raw grammar nor its resource grading is
reconstructed by \eqref{cat-eq:op}. Its interpreted image is the executable
part of the reconstructed operations, but that description is not a
recognition theorem. At the enriched level even the word ``image'' needs
care: the pointwise image of a functor need not be a subcategory when
distinct intermediate source objects become equal and new semantic
composites become possible. A full image would additionally insert
unrealized certificates. One can retain marked native operations as
objects with their actual native cells, or prove a suitable
composability-lifting condition; one may not assume either property from
an ordinary homwise set image.


\subsection{Existing generic arities already give recognition}
For a monad $T$ on $\Set$, with forgetful functor
$U:T\text{-}\mathsf{Alg}\to\Set$, the free algebra $F_Tn$ gives
\begin{equation}\label{eq:allmodels}
              \Nat(U^n,U)\cong Tn.
\end{equation}
Indeed, a natural operation is determined by
$t=\alpha_{F_Tn}(\eta_n)$, and naturality for the unique homomorphism
$F_Tn\to A$ extending $x:n\to UA$ gives
$\alpha_A(x)=a(Tx(t))$. Conversely that formula is natural. If $T$ is the
term monad of a finitary grammar, $t$ has a finite-term representative.
This is an actual recognition theorem, not the assertion that a supplied
semantic image equals itself. It uses the free arity model and all its
evaluation maps. It says nothing about a narrower resource fragment unless
that fragment is also recognized.

The closest general enriched result is Lucyshyn-Wright--Parker
Proposition~4.34: if $G:\mathcal A\to\mathcal C$ has the appropriate
$j$-relative left adjoint $F$, its structure is the opposite of the full
enriched subcategory on the free arity objects, with
\begin{equation}\label{eq:lwp-free}
 \Str(G)(J,K)\cong\mathcal A(FK,FJ)
                   \cong\mathcal C(jK,GFJ).
\end{equation}
Theorem~4.35 gives algebraicity criteria from these relative free objects,
their density and the specified nerve-image condition; strict algebraicity
also requires a discrete isofibration. Their base and arity hypotheses,
including the relevant enriched density and existence conditions, must be
checked \cite{LWP}. The reversed free-model hom in
\eqref{eq:lwp-free} is deliberate. An arrow $FK\to FJ$ is not thereby a
native execution in the reverse direction.

When free-model arrows are precisely the admitted native grammar, this
published result already discharges recognition. The problem pursued here
is the residual case where the test class contains only finite models,
where the generic free algebra is infinite, or where an independently
supplied native grammar is narrower than the ambient reconstructed theory.
For finite carrier clones, the classical $\operatorname{Pol}$--
$\operatorname{Inv}$ correspondence gives another genuine recognition
mechanism: a clone is recovered from the relations it preserves
\cite{Geiger}. It is not merely commutation with endomorphisms, need not
provide a small relation basis, and does not recognize an arbitrary
cost-bounded fragment unless its structural closure is justified.

\subsection{Typed change-of-doctrine data}
Let a second doctrine have data $\W,\E',\C',U',j'$ and $H'_{a'}$.
The following are sufficient, explicitly typed comparison data:
\begin{enumerate}[label=(\roman*)]
\item a strong symmetric monoidal functor $F:\V\to\W$, allowing change of base
      $F_*\E$;
\item a $\W$-functor $R:\E'\to F_*\E$;
\item a map of arity labels $\kappa:a\mapsto\kappa a$;
\item $\W$-natural isomorphisms
\begin{equation}\label{cat-eq:theta}
 \theta_a:H'_{\kappa a}\xrightarrow{\ \cong\ }
       (F_*H_a)R.
\end{equation}
\end{enumerate}
Here $F_*H_a$ denotes the $\W$-functor into the canonical
self-enrichment of $\W$ obtained from $H_a$ using the evaluation-induced
map $F[X,Y]\to[FX,FY]$. No preservation of internal homs is asserted.
The functor $R$ is a reduct or interpretation of \emph{target models}
in changed-base source semantics. It goes opposite to a source-program
compiler. In a structured comparison of arity categories, $\kappa$ and
the $\theta_a$ must also be natural in the arities; the labelwise statement
below only needs the displayed data. New target arities outside the image
of $\kappa$ are not covered by the comparison.

\begin{theorem}[Typed reconstruction comparison]\label{cat-thm:base}
The data above determine a $\W$-functor
\begin{equation}\label{cat-eq:comparison}
 \Phi_u:F_*\Op_d\longrightarrow \Op_{d'}|_\kappa.
\end{equation}
Here the target is reindexed to the source arity labels, with hom-object
$\Op_{d'}(\kappa a,\kappa b)$ at $(a,b)$. Its component at an ordinary operation is
\[
 (\Phi_u\alpha)_{E'}=
 \theta_b^{-1}\,(F\alpha)_{R E'}\,\theta_a.
\]
No preservation of the ends in \eqref{cat-eq:op} is needed to construct this
map, and the map need not be full, faithful, or essentially surjective.

For successive data $(F,R,\kappa,\theta)$ and
$(G,S,\lambda,\psi)$, their composite has base $GF$, reduct
$F_*\E\xleftarrow{R}\E'$ changed by $G$ and then precomposed with $S$,
that is, reduct $(G_*R)S$, arity map $\lambda\kappa$, and comparison
\[
 \theta^{\rm comp}_a=(G\theta_a)S\circ\psi_{\kappa a}.
\]
Under the canonical change-of-base identifications,
\begin{equation}\label{cat-eq:basecomposition}
 \Phi_{v\circ u}=\Phi_v\circ G_*\Phi_u.
\end{equation}
These comparisons satisfy unit and associativity coherence.
\end{theorem}
\begin{proof}
The enriched evaluation morphisms
$\Op_d(a,b)\otimes H_a(E)\to H_b(E)$ become, using strong
symmetric monoidality, morphisms
$F\Op_d(a,b)\otimes FH_a(E)\to FH_b(E)$.
Whiskering with $R$ and conjugating by $\theta$ gives a $\W$-natural
family with input $H'_{\kappa a}$ and output $H'_{\kappa b}$.
The universal property of the target enriched hom-object gives
$F\Op_d(a,b)\to\Op_{d'}(\kappa a,\kappa b)$.
Units and composition follow from those of evaluation and cancellation of
$\theta_b\theta_b^{-1}$. For two changes, substituting the displayed
composite $\theta$ yields exactly successive conjugation. For three
changes both parenthesizations give the same iterated conjugation, after
the monoidal coherence identifications. This proves coherence rather
than merely postulating a commuting square.
\end{proof}

\begin{remark}[A noninvertible arity comparison is a real obligation]
An arrow $H'_{\kappa a}\to(F_*H_a)R$ without an inverse does not supply
the output leg of \eqref{cat-eq:comparison}. Arrows in the reverse direction
do not supply its input leg. With quotient comparisons one may instead
restrict to operations that descend through their kernel relations and
prove descent; this is an additional theorem. Choosing unrelated lifts
at the input and output does not in general preserve identity or
composition. In particular, changing which diagrams count as inputs is
not just renaming the arity labels.
\end{remark}


\subsection{The native square and the claims it does not imply}
Let $c_u:\mathsf P_d\to\mathsf P_{d'}$ be an actual translation of
source primitives and structural rules. Correctness means, on its declared
operations and cells,
\begin{equation}\label{eq:compiler-square}
 e_{d'}c_u=\Phi_u e_d,
\end{equation}
with specified coherent isomorphisms in a weak setting. For a presentation by
generators and equations, correctness follows by checking primitive
interpretations and the defining equations, then inducting on expressions.
It is not supplied by an arbitrary functor between semantic carriers.

Equation~\eqref{eq:compiler-square} preserves native realizability.
Reflection would assert that a source semantic operation whose image has a
target native term itself has a source native term. Lifting would assert
that a target operation has a source preimage. These are different claims;
neither follows from the square. A right inverse on objects is weaker than
a natural, substitution-compatible section, and a section is weaker than
an equivalence retaining original executions.

A forward resource bound must also be proved from the compiler. If one
target constructor copies a translated input, its cost can grow with every
copy. Bounds compose by function composition, not by silently taking the
larger of two budgets. Reflection of a semantic operation need not recover
the original program, the cheapest program, its proof history, or its
physical runtime. The next two sections prove an extraction criterion with
explicit, independently checkable hypotheses.

\section{Finite grammars and independent arity probes}\label{alg-sec:finite-probes}
Fix a finite set of sorts $S$. Let $\mathcal V$ be a finitary many-sorted
variety over $\Set^S$, presented by a finite signature and equations.
All operations and equations are supplied, rather than inferred from a
finite transcript. Write $F_{\mathcal V}X$ for the free algebra on a finite
sorted set $X$. We use the usual set-based convention allowing empty sorts
when the signature permits them; alternatively restrict to arities for
which all free and generated algebras used below exist.

Let $\mathcal K$ be a full class of finite $\mathcal V$-algebras, closed
under finite products and subalgebras. It is sufficient to take a small
skeleton. A finite algebra has finitely many elements in total across its
sorts. Its underlying sorted set is $UA$. Define
\[
 \Op_{\mathcal K}(X,s)
 =\Nat\bigl(\Set^S(X,U-),U_s\bigr).
\]
Thus $\alpha\in\Op_{\mathcal K}(X,s)$ assigns a value
$\alpha_A(x)\in A_s$ to each finite model and sorted assignment
$x:X\to UA$, naturally in every admitted homomorphism.
A term $t\in(F_{\mathcal V}X)_s$ induces such an operation $[t]$.
The extensional native core is the image
\[
 \Core_{\mathcal K}(X,s)
 =\{[t]:t\in(F_{\mathcal V}X)_s\}.
\]
This definition only names the object to be recognized. The theorems below
give finite-test criteria that do not presuppose a native lift.
Native recipes remain marked separately if their identity, resource use,
or proof histories matter. Terms identified in every finite model need not
have the same native execution history.

With tuples of output terms, these cores form the usual many-sorted
operation theory under substitution. The full semantic operation families
also compose by substitution. Interpretation is a morphism of these
structures. No assertion is made that an arbitrary hom-category image
of a functor is a witness-preserving enriched subcategory.

\subsection{From finite interpolation to a recognition certificate}

\begin{lemma}[Finite simultaneous interpolation]\label{alg-lem:interpolation}
For every $\alpha\in\Op_{\mathcal K}(X,s)$ and every finite list of
pointed models $(A_i,x_i)$, there is a native term $t$ such that
\[
             \alpha_{A_i}(x_i)=t^{A_i}(x_i)
             \qquad\text{for every }i.
\]
\end{lemma}
\begin{proof}
Take the homomorphism
$F_{\mathcal V}X\to\prod_iA_i$ induced by the assignments and let $P$
be its image. It is finite and lies in $\mathcal K$. Let
$\bar x:X\to UP$ be the generating assignment. By definition of an
algebraic image, $\alpha_P(\bar x)$ is represented by some term $t$.
Naturality for $P\to A_i$ then proves all the equalities at once.
\end{proof}

Finite interpolation does not give one term working on all models.
For example, a sufficiently large power represents the idempotent-power
operation on any fixed finite list of finite monoids, but no fixed monoid
word represents it on every finite monoid. That is an obstruction for the
straight-line monoid grammar, not for every algorithmic language.

Give native syntax a cost with finitely many terms of cost at most $r$,
up to the supplied syntactic equations, and let $B_r$ be that finite set.
Finite signatures with a bound on expression size provide such budgets.

\begin{proposition}[The uniform bound is exactly what finite interpolation
lacks]\label{alg-prop:uniform}
An operation $\alpha$ has a representative in $B_r$ if and only if for
every finite pointed model $(A,x)$ there is some $t\in B_r$ with
$\alpha_A(x)=t^A(x)$. Equivalently, the same budget $r$ suffices for every
finite simultaneous list of pointed tests.
If the criterion fails, one finite pointed model refutes every term in
$B_r$ simultaneously.
\end{proposition}
\begin{proof}
The forward implication is immediate. If no term in $B_r$ represents
$\alpha$ globally, choose for each $t\in B_r$ a pointed finite model
$(A_t,x_t)$ separating it from $\alpha$. In the finite product, naturality
for projections shows that no member of $B_r$ matches $\alpha$.
One can instead use its generated subalgebra. This is a single finite
countertest. The argument for simultaneous lists is identical.
\end{proof}

This is a compactness argument for a finite resource budget, not an
automatic decision procedure. A positive finite certificate needs an
additional grammar property.

\begin{definition}[Finite arity recognizer]
A finite recognizer for a native term $t\in(F_{\mathcal V}X)_s$ is a
homomorphism $q:F_{\mathcal V}X\to Q$ with $Q\in\mathcal K$ and
\begin{equation}\label{alg-eq:singleton}
                    q_s^{-1}\{q_s(t)\}=\{t\}.
\end{equation}
The fibre equality is in the free algebra of the supplied grammar, not
merely among already known low-cost terms. It is a statement that can be
proved from a truncation construction, independently of $\alpha$.
\end{definition}

\begin{theorem}[One finite arity test recognizes a native operation]
\label{alg-thm:recognition}
Let $q$ recognize $t$, and let $\eta:X\to UF_{\mathcal V}X$ be the free
generators. Then for every natural operation $\alpha$,
\[
 \alpha=[t]
 \quad\Longleftrightarrow\quad
 \alpha_Q(Uq\,\eta)=q(t).
\]
\end{theorem}
\begin{proof}
Only the reverse implication needs proof. Fix $(A,x)$ and form the image
$P$ of the homomorphism $F_{\mathcal V}X\to Q\times A$ induced by
$Uq\,\eta$ and $x$. Since $P$ is a finite generated subalgebra it lies
in $\mathcal K$. Choose $u\in(F_{\mathcal V}X)_s$ representing
$\alpha_P(\bar x)$. Naturality for $P\to Q$ gives $q(u)=q(t)$, so
\eqref{alg-eq:singleton} gives $u=t$. Naturality for $P\to A$ now gives
$\alpha_A(x)=t^A(x)$. The pointed model was arbitrary.
\end{proof}

\begin{corollary}[A complete bounded recognizer]\label{alg-cor:budget}
Suppose a finite algebra $T_r\in\mathcal K$ and map $q_r:F_{\mathcal V}X\to T_r$
have singleton fibres at the images of every term in $B_r$, and those
images are known and distinct. Then
\[
 \alpha\in[B_r]
 \quad\Longleftrightarrow\quad
 \alpha_{T_r}(Uq_r\eta)\in q_r(B_r).
\]
The value on this one finite arity test directly identifies a native term
of budget at most $r$.
\end{corollary}

\subsection{Independent constructions of the recognizers}
For an absolutely free finite typed tree grammar, keep all terms of height
at most $r$ and add a typed overflow value. A constructor returns its
actual tree if the height remains at most $r$, and overflow otherwise;
any overflow input produces overflow. This is a finite algebra, and every
non-overflow fibre is a singleton. The essential property is that an
oversized subtree cannot be erased by an enclosing native constructor.

For a free monoid on finite $X$, use
\[
 T_r=X^{\leq r}\sqcup\{\bot\},
 \qquad
 u\cdot v=
 \begin{cases}uv&|uv|\leq r,\\ \bot&|uv|>r,\end{cases}
\]
with absorbing $\bot$. Its size is $1+\sum_{j=0}^r|X|^j$.
For a free category on a finite typed graph, use the analogous typed path
truncation, with an absorbing zero arrow for each ordered pair of objects.
Every path of length at most $r$ is retained as a distinct arrow. The
quotient need not hit a zero arrow in a hom-set with no sufficiently long
path; surjectivity of $q_r$ was not required.

A useful general sufficient condition is a finite low-cost part together
with tails forming a typed ideal under all contexts, so that collapsing
the tails is a congruence. This condition fails for grammars with genuine
cancellation or erasing contexts. Residual finiteness alone is not enough:
it separates any two fixed terms, whereas a recognizer must separate a
term from all other terms at once. An infinite free group has no
singleton fibre under a homomorphism to a finite group.

\begin{proposition}[An intrinsic topological description under the
recognizer hypothesis]\label{alg-prop:isolated}
Give $\Op_{\mathcal K}(X,s)$ the topology of pointwise evaluation on
finite pointed models, each response set being discrete. If every native
term has a finite recognizer, then
\[
       \Core_{\mathcal K}(X,s)
       =\{\text{isolated points of }\Op_{\mathcal K}(X,s)\}.
\]
\end{proposition}
\begin{proof}
A recognizer makes the operation's singleton an open evaluation fibre by
Theorem~\ref{alg-thm:recognition}. Conversely, a basic open neighborhood of an
isolated operation fixes finitely many pointed evaluations.
Lemma~\ref{alg-lem:interpolation} supplies a native term in that neighborhood,
which must therefore be the operation itself.
\end{proof}
The topological description is conditional on the supplied grammar's
finite recognizability. It is not a universal way to reconstruct an
unspecified execution repertoire from semantics. One may add an implicit
operation as a new primitive while leaving the same finite semantic
models and their underlying functor unchanged.

\section{A native-reflection theorem under a change of doctrine}

Let $\mathcal W$ be another finitary many-sorted variety on the same
finite sort set. A primitive term interpretation
$\sigma:\mathcal V\to\mathcal W$ replaces each source primitive by a
finite target term, and sends every source equation to a valid target
equation. It induces a reduct functor
\[
           R:\mathcal W\longrightarrow\mathcal V
\]
with unchanged underlying sorted carriers. Let $\mathcal L$ be a full
class of finite $\mathcal W$-algebras closed under products and subalgebras,
and assume $R(\mathcal L)\subseteq\mathcal K$.

There is a well-typed semantic restriction
\begin{equation}\label{alg-eq:comparison}
 C_R:\Op_{\mathcal K}(X,s)\longrightarrow\Op_{\mathcal L}(X,s),
 \qquad (C_R\alpha)_B(x)=\alpha_{RB}(x).
\end{equation}
For a source term, $C_R[t]=[\sigma t]$. Hence this map preserves the native
core, the structural operations, and substitution. This is the
straightforward preservation statement; reflection requires more.

\begin{definition}[Finite source-test coverage]\label{alg-def:cover}
Say that $R$ covers the finite source tests if every $A\in\mathcal K$
is a homomorphic image of a subalgebra of a finite product
\[
                   \prod_{i=1}^m RB_i,
                   \qquad B_i\in\mathcal L.
\]
All covers are ordinary source-algebra maps. This is a finite HSP
condition on the model classes. It does not assume anything about the
realizability of an unknown natural operation.
\end{definition}

\begin{lemma}[Finite source-test coverage reflects semantic equality]
\label{alg-lem:faithful}
Under Definition~\ref{alg-def:cover}, $C_R$ is injective on every operation
family $\Op_{\mathcal K}(X,s)$.
\end{lemma}
\begin{proof}
Suppose $C_R\alpha=C_R\beta$. On a finite product of reducts, the
projections and naturality show $\alpha=\beta$ pointwise. For a subalgebra
$D$ of that product, the injective inclusion and naturality give equality
on $D$. For a surjection $e:D\twoheadrightarrow A$, lift any assignment
$x:X\to UA$ to $\bar x:X\to UD$. Naturality gives
\[
 \alpha_A(x)=e\alpha_D(\bar x)=e\beta_D(\bar x)=\beta_A(x).
\]
Every finite source model has such a cover.
\end{proof}

\begin{theorem}[Primitive compiler-image recognition]\label{alg-thm:compilerimage}
Let $t\in(F_{\mathcal W}X)_s$ have a finite recognizer
$q:F_{\mathcal W}X\to Q$ with $Q\in\mathcal L$ (and hence $RQ\in\mathcal K$). Write
$x=Uq\eta$ and let $D\subseteq RQ$ be the source subalgebra generated by
$x$. Then the following are equivalent:
\begin{enumerate}
\item There is a source term $u$ with $\sigma u=t$ in the free target
      algebra $F_{\mathcal W}X$.
\item $q(t)\in D_s$.
\item The target natural operation $[t]$ is the restriction of a source
      term operation.
\end{enumerate}
These conditions are decidable from the finite tables of $Q$, the
primitive compiler, the marked generator assignment $x$, and $q(t)$. When they hold, finite source-subalgebra
closure extracts a witnessing $u$. Moreover, if $[t]$ extends to any
source natural operation at all, then these equivalent conditions hold.
No finite source-test coverage assumption is needed for this theorem.
\end{theorem}
\begin{proof}
A source term evaluates into $D$, proving (1)$\Rightarrow$(2).
If (2) holds, generation supplies a source term $u$ with
$u^{RQ}(x)=q(t)$. By the primitive interpretation,
$q(\sigma u)=u^{RQ}(x)=q(t)$, and the singleton fibre of $q$ at $t$
forces $\sigma u=t$. Thus (2)$\Rightarrow$(1), and (1)$\Rightarrow$(3)
is immediate. For (3)$\Rightarrow$(2), evaluate the restricting source
term at $RQ,x$.

More generally, for an extension $\alpha$, naturality for the inclusion
$D\hookrightarrow RQ$ gives
$q(t)=\alpha_{RQ}(x)\in D$. Finite subalgebra closure is an effective
procedure on finite tables, retaining a term for every generated element.
It either finds $q(t)$ or proves that no source term can compile to $t$.
\end{proof}

This is an independently testable source-realization criterion: it
reduces a question in an infinite source grammar to generation inside a
specific finite algebra computed from the target recognizer. The
singleton-fibre hypothesis, rather than residual finiteness alone, is
what makes this finite negative answer sound.

\begin{theorem}[Conservative native transport with extraction]
\label{alg-thm:transport}
Assume finite source-test coverage. Assume also that every target native
term has a finite recognizer in $\mathcal L$. Then for every source
natural operation $\alpha$,
\begin{equation}\label{alg-eq:reflectnative}
 \alpha\in\Core_{\mathcal K}(X,s)
 \quad\Longleftrightarrow\quad
 C_R\alpha\in\Core_{\mathcal L}(X,s).
\end{equation}
More precisely, if $C_R\alpha=[t]$ and a finite recognizer
$q:F_{\mathcal W}X\to Q$ for $t$ is supplied, a source term realizing
$\alpha$ can be constructed from the finite reduct $RQ$ and the value
$q(t)$. No source decompiler is assumed. The output is a term modulo the supplied
native equations, with an explicit raw expression retained by the algorithm.

If the finite tables and primitive interpretation are effective, construct
the source subalgebra of $RQ$ generated by $Uq\eta$ while retaining an
expression for each new element, and return any retained expression for
$q(t)$. This procedure terminates. If $M=\sum_{a\in S}|(RQ)_a|$, the
returned source term can have depth at most $M$. For maximum primitive
arity $d\geq2$, one may bound its syntax-tree size by
$(d^{M+1}-1)/(d-1)$; for $d=1$ the bound is $M+1$, and for $d=0$ a returned expression
is a variable or a constant. These are syntax bounds, not optimal runtime
or physical-resource bounds.
\end{theorem}
\begin{proof}
Preservation follows from primitive compilation. For reflection, put
$x=Uq\eta$ and let $D$ be the source subalgebra of $RQ$ generated by $x$.
The inclusion $i:D\hookrightarrow RQ$ and naturality imply
\[
 q(t)=(C_R\alpha)_Q(x)=\alpha_{RQ}(x)
       =i\alpha_D(\bar x)\in i(D).
\]
Consequently there is a source term $u$ such that $u^{RQ}(x)=q(t)$.
Its compiled target term satisfies
\[
          (\sigma u)^Q(x)=u^{RQ}(x)=q(t).
\]
Theorem~\ref{alg-thm:recognition}, applied to the target recognizer, shows
$[\sigma u]=[t]$ on every model in $\mathcal L$. Thus
$C_R[u]=C_R\alpha$, and Lemma~\ref{alg-lem:faithful} gives $[u]=\alpha$ on
all of $\mathcal K$.

For the algorithm, initialize with the variables and nullary primitive
values and iteratively apply every primitive to known inputs. Because
there are at most $M$ elements altogether, every strictly growing round
adds an element and the process stabilizes after at most $M$ rounds.
The displayed membership guarantees that $q(t)$ is reached. An expression
produced by round $k$ has height at most $k$ with variables and constants
at height zero. The tree-size bound follows from the geometric sum for
branching at most $d$. No effective presentation or equality procedure for
the infinite free algebra is needed during this finite extraction.
\end{proof}

\paragraph{What is genuinely being checked.}
The finite recognizer is constructed from the target free grammar. The
coverage theorem is proved using finite source models, not by assuming
that $C_R$ already reflects executable operations. The extraction uses
only a finite target test and the source primitives. The hypotheses can
fail independently, and neither follows from a formal change-of-base
construction alone.

\paragraph{What the theorem does not say.}
The two sides of \eqref{alg-eq:reflectnative} concern semantic operations
already related by $C_R$. The theorem does not say that every target
native operation extends to a source natural operation. If such an
extension does exist, it supplies a source implementation. The equality
$C_R\alpha=[t]$ is an actual premise about the semantic operation, not
something established by evaluating an arbitrary unknown procedure once.
The construction transports a source equation only relative to the
specified finite model class; it does not produce an intensional proof
history or identify a physical implementation without its own warrant.

\subsection{Composition of doctrine changes}

\begin{theorem}[Coherent composition and transitivity of coverage]
\label{alg-thm:composition}
For primitive interpretations
$\mathcal V\xrightarrow{\sigma}\mathcal W
\xrightarrow{\tau}\mathcal Z$ with reducts $R$ and $S$, the semantic
comparisons and native compilation satisfy
\[
                     C_S C_R=C_{RS},
             \qquad \tau(\sigma t)=(\tau\sigma)t.
\]
If $R$ and $S$ each cover the corresponding finite source tests, then
$RS$ does too. Hence the native-reflection theorem applies directly to
the composite whenever the final target has finite recognizers.
Whenever the recognizers required at both stages are available,
successive extraction and direct extraction give the same extensional
source operation. They need not give the same syntactic recipe or native
certificate.
\end{theorem}
\begin{proof}
Both comparison formulas are equalities on each model and assignment.
For transitivity, cover a source algebra $A$ by a subalgebra
$D\subseteq\prod_i RB_i$ and a surjection $D\twoheadrightarrow A$.
Cover each $B_i$ by $E_i\subseteq\prod_j SC_{ij}$ with a surjection
$E_i\twoheadrightarrow B_i$. Reducts induced by term interpretations
preserve products, subalgebras, and underlying surjections. Thus
$\prod_i RE_i\to\prod_i RB_i$ is surjective and its domain is a
subalgebra of a finite product of $RSC_{ij}$.
The inverse image of $D$ is a subalgebra still surjecting onto $D$, and
then onto $A$. This is the required composite cover. The extraction
claim follows because every extracted term realizes the same $\alpha$.
\end{proof}

\subsection{A nontrivial positive change}
Take source semigroups and target monoids, with reduct forgetting the
unit, and a nonempty finite variable set. Every finite semigroup embeds
in its finite unitization, so finite source-test coverage holds.
Finite monoid word truncations provide the target recognizers.
Consequently a natural operation on finite semigroups is a nonempty
finite word exactly when its restriction to finite monoids is a finite
monoid word.

The empty target word does not secretly create a source operation.
In a fresh unitization of a finite semigroup, the subsemigroup generated
by nonempty-variable inputs stays in the original semigroup and misses
the newly adjoined unit. The extraction proof therefore shows that the
constant new unit cannot be the restriction of such a source natural
operation. The theorem reflects an existing operation, rather than
assuming that every target operation has an extension.

This example changes the algebraic operation doctrine but not enrichment.
A genuinely witness-forgetting change is considered in
Section~\ref{sec:thinning}, where the positive reflection hypotheses fail
in an informative way.


\section{Finite monoids as a complete algebraic instance}
\label{sec:monoids}
This section makes the preceding framework explicit on an entire category
of models, not on one fixed monoid. The primitive grammar is ordinary
monoid terms in $1$ and multiplication, with cartesian variable use.
Let $U:\FinMon\to\Set$ be the underlying functor and let
$X=\{x_1,\ldots,x_n\}$ with $n\geq1$. Native terms modulo the monoid
axioms are words $X^*$, and their cost is word length, counting variable
occurrences. Semantic operations are
$\Op_U(n)=\Nat(U^n,U)$. Thus, for every homomorphism $h:M\to N$,
\begin{equation}\label{mon-eq:nat}
 h(\alpha_M(\mathbf a))=\alpha_N(h\mathbf a).
\end{equation}
A word $w$ acts by evaluation $\ev_w=[w]$, and write
$e_n:X^*\to\Op_U(n)$ for $w\mapsto\ev_w$.
The components of a semantic operation need not be monoid homomorphisms.
Multiplication is a natural operation although $M^2\to M$ need not be a
homomorphism for a noncommutative $M$.

All monoids, rather than only finite ones, include the free monoid $X^*$,
so \eqref{eq:allmodels} already recovers exactly the words. Restricting to
finite monoids changes that answer as follows.

\subsection{Semantic completion by profinite words}
\label{mon-sec:profinite}
Consider all finite quotient homomorphisms $q:X^*\twoheadrightarrow Q$,
with connecting homomorphisms $h:Q\to Q'$ satisfying $hq=q'$. Their inverse
limit is the free profinite monoid $\widehat{X^*}$, equipped with the
subspace topology from the product of its finite discrete quotients.
Forgetting topology gives the following standard correspondence.

\begin{proposition}[Implicit operations and profinite words]
\label{mon-prop:implicit}
There is a canonical bijection
\[
 \Op_U(n)\cong\bigl|\widehat{X^*}\bigr|
\]
under which $e_n$ is the ordinary-word inclusion into its profinite
completion.
\end{proposition}
\begin{proof}
From $\alpha$ obtain the coordinate
$a_q=\alpha_Q(q(x_1),\ldots,q(x_n))$. Equation~\eqref{mon-eq:nat} makes
these coordinates compatible with connecting maps.

Conversely, let $(a_q)_q$ be compatible. For $\mathbf b\in M^n$, let
$S\leq M$ be the submonoid generated by the entries of $\mathbf b$ and
let $q_{\mathbf b}:X^*\twoheadrightarrow S$ evaluate the variables there.
Set $\alpha_M(\mathbf b)=i_S(a_{q_{\mathbf b}})$, where $i_S:S\hookrightarrow M$.
For $h:M\to N$, the restriction $h|_S$ is a quotient onto the submonoid
generated by $h\mathbf b$. Compatibility therefore gives
\eqref{mon-eq:nat}. The two constructions are inverse: for a natural operation,
naturality of $i_S$ gives the displayed formula, and for a quotient $q$ its
generator tuple generates all of $Q$. For a word $w$, its coordinate is
$q(w)$, exactly the usual comparison with the completion.
\end{proof}

This is not a new representation theorem. In Ad\'amek--Chen--Milius--Urbat,
Definition 5.16 calls these natural transformations \emph{implicit
operations}; Lemma 5.19 identifies them with profinite terms
\cite{ACMU}. Their Definition 4.5 constructs the profinite monad as a
codensity monad on $\mathbf{Stone}$, and Construction 4.6 gives its limit
formula. At finite discrete variable sets its underlying set is exactly
that in Proposition~\ref{mon-prop:implicit}. The Set-valued right Kan extension
$\operatorname{Ran}_U U$ has
$(\operatorname{Ran}_U U)(X)=\Nat(\Set(X,U-),U)$; at finite $X$ this is
our same set. These are related constructions on different base categories,
not an identification of their entire monads without the required change
of base. The distinction is expressly discussed in their Remark 4.7.

\subsection{The idempotent power is computable but is not one word}
\label{mon-sec:omega}
\begin{lemma}[Uniform exponent at a finite size bound]
\label{mon-lem:factorial}
If $M$ is a finite monoid of size $m\geq1$ and $a\in M$, then
$a^{k!}$ is the same idempotent for every integer $k\geq m$.
It is the unique idempotent among the positive powers of $a$.
\end{lemma}
\begin{proof}
Among $a,a^2,\ldots,a^{m+1}$ two powers coincide. Hence there are
$r,p\geq1$ with $r\leq m$, $p\leq m$, and
$a^{j+p}=a^j$ for every $j\geq r$. All exponents at least $r$ that
are divisible by $p$ therefore give the same element. For $k\geq m$,
we have $k!\geq r$ and $p\mid k!$, including $m=1$ directly.
Since both $k!$ and $2k!$ satisfy these conditions, $a^{k!}$ is idempotent,
and all the asserted factorial powers agree. Finally, if a positive power
$a^s$ is idempotent, then $a^{js}=a^s$ for every $j\geq1$. Choose $j$
so that $js\geq r$ and $p\mid js$. This gives
$a^s=a^{js}=a^{k!}$ and proves uniqueness without assuming that the
originally selected period $p$ was minimal.
\end{proof}

For $m=|M|$, define $\omega_M(a)=a^{m!}$, also written $a^\omega$.
For a homomorphism $h:M\to N$, choose $k\geq\max(|M|,|N|)$;
then $h(a^\omega)=h(a^{k!})=h(a)^{k!}=h(a)^\omega$.
Thus $\omega\in\Op_U(1)$. This is precisely the standard example in
\cite[Example 5.17]{ACMU}.

The family is uniformly computable from finite multiplication tables.
One may compute $a^{m!}$ by repeated squaring, or enumerate positive powers
until the first repetition and take the power in the eventual cycle whose
exponent is divisible by its period. This is a terminating finite-table
algorithm; no claim about its succinct-input complexity is needed here.
For each fixed $M$ the single word $x^{m!}$ represents $\omega_M$ at all
inputs. The truncation calculation below proves that no fixed word represents
the whole family.

\subsection{The finite probe and its exact resource threshold}
\label{mon-sec:probe}
For $L\in\N$, define the finite monoid
\[
 Q_{X,L}=X^{\leq L}\sqcup\{\bot\}.
\]
Its multiplication is concatenation if the result has length at most $L$,
and $\bot$ otherwise; $\bot$ is absorbing and $\eps$ is the identity.
It is the Rees quotient of $X^*$ by the ideal of words longer than $L$.
Associativity follows either from that quotient or from the observation
that any product whose total length exceeds $L$ overflows regardless of
parenthesization. Write $q_L:X^*\twoheadrightarrow Q_{X,L}$ and
$\mathbf x_L=(q_L(x_1),\ldots,q_L(x_n))$. When $L=0$, each variable is
sent to $\bot$; this is deliberate and the results include that case.

\begin{theorem}[Truncation certificate and exact native extraction]
\label{mon-thm:probe}
For $\alpha\in\Op_U(n)$ and $L\geq0$, the following are equivalent:
\begin{enumerate}[label=(\alph*)]
\item $\alpha=\ev_w$ for some $w\in X^{\leq L}$;
\item $\alpha_{Q_{X,L}}(\mathbf x_L)\ne\bot$;
\item for every finite $M$ and every $\mathbf a\in M^n$ there is a word
      $v_{M,\mathbf a}\in X^{\leq L}$, allowed to depend on both,
      with $\alpha_M(\mathbf a)=v_{M,\mathbf a}^M(\mathbf a)$.
\end{enumerate}
If (b) returns the word $w$, then this very word is the unique native
representative of $\alpha$. Consequently $\alpha$ is native if and only
if (b) holds for some $L$.
\end{theorem}
\begin{proof}
The implications (a)$\Rightarrow$(c)$\Rightarrow$(b) follow by evaluation,
since a word of length at most $L$ does not overflow on $\mathbf x_L$.
For completeness, specialize the proof of Theorem~\ref{alg-thm:recognition}.
For (b)$\Rightarrow$(a), fix an arbitrary finite $M$ and tuple $\mathbf a$.
Inside $Q_{X,L}\times M$ take the finite submonoid
\[
 P=\langle(q_L(x_i),a_i):1\leq i\leq n\rangle.
\]
Write $\mathbf b$ for its specified generating tuple. Because $P$ is
generated by $\mathbf b$, the element $\alpha_P(\mathbf b)$ equals
$v^P(\mathbf b)$ for some finite word $v$. Naturality of the first
projection gives
\[
 q_L(v)=\alpha_{Q_{X,L}}(\mathbf x_L)=w.
\]
The fiber $q_L^{-1}(w)$ is the singleton $\{w\}$, so $v=w$.
Naturality of the second projection now gives
$\alpha_M(\mathbf a)=w^M(\mathbf a)$. This proves the global statement
without assuming that a finite free monoid exists. Finally, distinct
words $w,v$ are distinguished by $Q_{X,L'}$ at its generators whenever
$L'\geq\max(|w|,|v|)$, so native representatives are unique.
\end{proof}

The generated submonoid in this proof is essential. An arbitrary element
of the full product $Q_{X,L}\times M$ need not be the value of the same
word in its two coordinates. Using the full product without the generated
submonoid would leave precisely the reconstruction step unproved.

\begin{corollary}[The exact length profile]
\label{mon-cor:length}
If $\alpha=\ev_w$ and $|w|=k$, then its probe output is $\bot$ for
$L<k$ and is $w$ for every $L\geq k$. If $\alpha$ is nonnative, every
probe output is $\bot$. For unary $\omega$, every output is $\bot$,
so $\omega$ is not a native monoid term.
\end{corollary}
\begin{proof}
The first assertion is direct evaluation. A nonoverflow value would imply
nativeness by Theorem~\ref{mon-thm:probe}. In the unary truncation, every
positive sufficiently long power of $q_L(x)$ is $\bot$, including the
case $L=0$, so its idempotent power is $\bot$.
\end{proof}

\begin{remark}[Nullary operations]
At $n=0$ there is exactly one natural operation, the identity element.
The unique homomorphism from the one-element monoid to each $M$ forces it.
The free monoid on no variables is already the one-element monoid.
There is no additional nullary profinite element. The artificial overflow
point in a formally defined $Q_{\varnothing,L}$ is not generated by the
empty tuple, so treating it as a quotient of $\varnothing^*$ would be
incorrect; the nullary case is handled separately as above.
\end{remark}

\subsection{What the constructive statement actually supplies}
Suppose the input is an exact evaluation procedure for $\alpha$, together
with a valid promise or proof of naturality on \emph{all} finite monoids.
Given $L$, build the finite table of $Q_{X,L}$ and query
$\alpha_{Q_{X,L}}(\mathbf x_L)$. A nonoverflow answer returns a word and
Theorem~\ref{mon-thm:probe} certifies its meaning on every finite model. An
overflow answer certifies that no word of length at most $L$ represents
$\alpha$; it does not certify that no longer word does.

Increasing $L=0,1,\ldots$ is a positive semidecision procedure. It halts
exactly on native operations. With a supplied bound on native length it is
a terminating reconstruction method; under no bound it need not halt.
The probe has
\[
 |Q_{X,L}|=1+\sum_{j=0}^{L}n^j
\]
elements. Thus one query can require an exponentially large explicit
multiplication table when $n>1$. Neither naturality nor physical access to
that probe is established merely by writing this mathematical procedure.
If an experiment supplies only a scalar readout rather than the full
output element, a justified separating readout family is additionally
needed to recover the value used in the theorem. No such access is inferred
from an ordinary response transcript.


\paragraph{Classical status and substitution.}
Finite words are exactly the isolated points of the free profinite monoid;
this is stated, for example, by Costa in Section~2.3 \cite{Costa}.
Together with the implicit-operation correspondence, this already entails
finite positive certificates for words. The truncation calculation supplies
the explicit probe, returned word and resource threshold for our present
question, not a new isolation theorem.

Native and semantic operations both form clones under substitution, and
word evaluation satisfies
\[
 [w[v_1/x_1,\ldots,v_k/x_k]]
       =[w][[v_1],\ldots,[v_k]],\qquad
 |w[v_1/x_1,\ldots,v_k/x_k]|
       =\sum_i\#_{x_i}(w)|v_i|.
\]
Because a native representative is unique in this example, successful
extraction respects substitution extensionally and returns its resulting
word. It need not preserve a naively chosen budget. For $K\geq L$, the
maps $Q_{X,K}\to Q_{X,L}$ that retain short words and overflow the rest
are homomorphisms, commute with the quotient maps and compose strictly.
Naturality makes the probe answers compatible with them.

\subsection{A doctrine change that fails reflection}
Let $J:\FinGrp\hookrightarrow\FinMon$ be the full inclusion.
Its reduct comparison is restriction and satisfies
\[
                 J^*(x^\omega)=1=J^*(1).
\]
It therefore reflects neither equality nor native monoid-term definability.
The source-test coverage hypothesis has genuinely failed: products,
submonoids and quotients of finite groups are finite groups, whereas a
finite monoid with a nontrivial absorbing element need not be a group.
The isolating truncation tests have been removed. Fullness and faithfulness
of the model inclusion do not repair the missing test coverage.

On the same category $\FinGrp$ with the same underlying functor, the
palettes $(1,\cdot)$ and $(1,\cdot,(-)^{-1})$ have different native cores.
Inversion is native only in the latter: a unary monoid term is $x^k$ for
$k\geq0$, and a cyclic group of order $k+2$ separates it from $x^{-1}$.
Thus the tested semantics cannot identify which of these two palettes was
authorized. This compares selected grammars on a finite model class, not
equivalent categories of \emph{all} models of two inequivalent algebraic
theories.

Likewise one may expand every finite monoid by its canonical $\omega$
operation without changing the category of finite models over $\Set$:
all homomorphisms already preserve it. The expanded grammar has one more
primitive. Naming it does not supply its execution resources or an
ordinary finitary equational presentation of this selected expanded class.


\section{Changing from categorical witnesses to reachability}
\label{sec:thinning}
The algebraic theorem changes a primitive grammar while retaining
set-valued carriers. A second test is needed when the doctrine itself
forgets witness identity. We now use exactly the comparison of
Theorem~\ref{cat-thm:base}, with $\rho:\Cat\to\Preord$ and the same
object-level presentations. This is not another name for profinite
completion: it can fail even when every category and every relevant
constraint is finite.

Let $\rho:\Cat\to\Preord$ send a small category to its preorder with the
\emph{same objects} and $x\le y$ iff some arrow $x\to y$ exists.
Do not silently pass to a skeletal poset: mutually reachable objects
remain distinct. The functor preserves products, so it is strong
cartesian monoidal. For a Cat-enriched presentation category $\E$ and
$U:\E\to\Cat$, put $\E'=\rho_*\E$ and let $U'$ be obtained by applying
$\rho$. For a finite discrete arity $n$,
\[
 \rho\Cat(n,U E)=\rho((U E)^n)\cong(\rho U E)^n
      =\Preord(n,\rho U E).
\]
Thus Theorem~\ref{cat-thm:base} applies with the identity reduct
$R:\E'\to\rho_*\E$ and the displayed arity isomorphisms. It does
\emph{not} say that $\rho$ preserves the reconstruction ends.

\begin{proposition}[Arrow-shaped arities do not transport automatically]
For the walking-arrow category $[1]$, the canonical map
\[
 \rho\Cat([1],C)\longrightarrow\Preord(\rho[1],\rho C)
\]
need not be an isomorphism, or even an equivalence of preorders.
\end{proposition}
\begin{proof}
Let $C$ have objects $0,1$, two distinct arrows $r,s:0\to1$, and only
identity arrows otherwise. The objects of $\Cat([1],C)$ are
$\id_0,\id_1,r,s$. In its reachability preorder, $r,s$ are incomparable:
a natural transformation between them would have identity components
and hence require $r=s$. The target has just the three monotone maps
$(0,0),(0,1),(1,1)$, with $r,s$ both sent to $(0,1)$.
These two incomparable source objects cannot become isomorphic through
an equivalence. This is loss of which witness was supplied as an input,
not merely a change of variable names.
\end{proof}

Let $W$ have two objects $a,b$, identities, and exactly two other arrows
$p,q:a\to b$. Let $G=C_2$ act by the automorphism $\sigma$ exchanging
$p$ and $q$ while fixing both objects. Regard $BG$ as a locally discrete
$\mathsf{Cat}$-enriched category and put $U:BG\to\mathsf{Cat}$ equal to
this action. Use finite discrete arities, so reconstructed binary
operations are equivariant functors $W^2\to W$, and their morphisms are
equivariant natural transformations. After reachability, $\rho W$ is the
chain $a<b$ and the $G$-action is trivial.

\begin{proposition}[A finite natural operation with no coherent lift]
\label{prop:dc-meet}
The monotone operation
$m:(\rho W)^2\to \rho W$, $m(x,y)=\min(x,y)$, is a natural binary operation of
the changed presentation functor. It is not the image of any natural
binary operation of $U$. In particular, a valid cartesian doctrine
change need not make the comparison from changed reconstructed structure
to reconstructed changed semantics surjective on operations.
\end{proposition}
\begin{proof}
Monotonicity and naturality of $m$ are immediate, the changed action
being trivial. Suppose an equivariant functor $F:W^2\to W$ lifts $m$.
For any $u,v\in\{p,q\}$, consider the commuting product square from
$(a,a)$ to $(b,b)$ whose first-coordinate edge is $u$ and whose
second-coordinate edge is $v$. The two initial legs end at $(b,a)$ and
$(a,b)$; all three of these objects have image $a$. Their image arrows
are therefore $1_a$. Functoriality of the square forces
\[
 F(1_b,v)=F(u,1_b) \quad\text{for every }u,v\in\{p,q\}.
\]
All these crossing arrows must consequently have one common image
$r\in\{p,q\}$. Equivariance applied to $(1_b,p)$ requires
\[
 r=F(1_b,q)=F(\sigma(1_b,p))=\sigma F(1_b,p)=\sigma r,
\]
which is impossible. Without the equivariance requirement there are
exactly two lifts, choosing all crossing arrows to be $p$ or all to be
$q$. Thus the obstruction is uniformity under an admitted presentation
map, not empty output hom-sets.
\end{proof}

The same example isolates witness loss in nullary operations. The
category of natural operations $1\Rightarrow U$ is the fixed-point
category $W^G$, which is discrete on $a,b$ because neither $p$ nor $q$
is fixed. But after $\rho$, the action is trivial, so the corresponding
operation preorder is  $\rho W$. Therefore
\[
 \rho(W^G)\longrightarrow (\rho W)^G
\]
is the discrete-two-points to two-element-chain map, not an isomorphism.
This is the concrete equalizer/end obstruction behind failure of
reconstruction to commute with change of enrichment. Product preservation
alone does not preserve the naturality constraints used in reconstruction.

Individual witness choices can also fail to compose, independently of
presentation symmetry. The exact path criterion below isolates that second
obstruction while retaining genuinely noninvertible arrows.

If native arrows carry path costs, reachability forgets those as well.
For example, declaring $p$ to have cost one and $q$ cost two sends two
witnesses of different resource requirements to the same inequality.
One must specify whether a budget asks for \emph{some} witness below the
bound, the particular retained witness, or \emph{all} witnesses. Bare
reachability preserves none of these distinctions. Adding minimum cost
can recover the first question on suitable finite presentations, but
cannot recover the latter two in general.

\subsection{A finite recognition theorem with a stated native machine model}
\label{subsec:dc-finite-recognition}
The preceding finite example permits a constructive recognition result
that does not assume lifts already exist. Its scope must be explicit:
a native operation is a finite functor table whose output arrows are
composites of admitted native generators. This is not automatically a
term in an arbitrary external programming language.

Let $\mathcal E$ be a finite category and let
$U:\mathcal E\to\mathsf{Cat}$ have finite category values $C_e$, supplied
by object, arrow and composition tables. In each $C_e$ specify finitely
many admitted native generating arrows with nonnegative integer costs;
identities cost zero. Let $N_e$ be the subcategory generated by them,
and let $d_e(h)$ be the minimum cost of an admitted path realizing $h$,
or infinity if none exists. These values are computable by finite
shortest-path closure on the supplied composition tables. A minimum is
attained whenever it is finite. No claim of polynomial complexity is
needed.

Suppose given a natural family of monotone object operations
$f_e:(\rho C_e)^n\to \rho C_e$ and an integer output budget $b_{e,\mathbf h}$
for each arrow tuple $\mathbf h$ of $C_e^n$. Form the following finite
constraint problem. For each $e$ and $\mathbf h:\mathbf x\to\mathbf y$
choose a variable
\[
 z_{e,\mathbf h}\in
 \{h\in C_e(f_e\mathbf x,f_e\mathbf y):
                  d_e(h)\leq b_{e,\mathbf h}\}.
\]
Require
\begin{align*}
 z_{e,1_{\mathbf x}}&=1_{f_e\mathbf x},\\
 z_{e,\mathbf k\mathbf h}&=z_{e,\mathbf k}z_{e,\mathbf h},\\
 U(v)(z_{e,\mathbf h})&=z_{e',U(v)^n(\mathbf h)}
       &&(v:e\to e').
\end{align*}
The given target naturality already supplies the corresponding object
equalities in the third equation.

\begin{proposition}[Finite source-table recognition]
\label{prop:dc-finite-csp}
The family $f$ has a natural $\mathsf{Cat}$-valued lift, executable in the
specified finite-table/native-generator model and meeting every displayed
budget, if and only if this finite constraint problem is satisfiable.
When it is, a lift and concrete native output paths can be extracted.
When it is not, unsatisfiability is an obstruction to this native lift,
not merely a failure of a chosen heuristic.
\end{proposition}
\begin{proof}
A lift supplies the variables by its arrow images. Functoriality,
naturality and execution bounds give all the constraints. Conversely,
a satisfying assignment, together with the given object functions,
defines functors $F_e:C_e^n\to C_e$ by the first two equations; the third
makes them a natural family. For each selected output arrow, retrieve a
minimum-cost native path from the finite closure calculation. The finite
tables and these paths supply the stipulated execution object. Enumeration
of the finite domains decides satisfiability, although its worst-case
complexity can be large. This is an explicit arity-level criterion from
primitive category tables and native generators; it does not quantify
over an uncomputed semantic image.
\end{proof}

Applied to $W$ and $BG$, the variables for the crossing arrows have
domains $\{p,q\}$. The composition equations identify them, and the
naturality equation forces the common value to be fixed by the swap.
This gives a transparent finite unsatisfiability certificate.
The theorem also exposes exactly what an additional native programming
language would require: a verified compiler from these finite tables and
paths, or a further grammar-recognition condition. That extra compiler
must not be inferred merely from the existence of the categorical lift.

\subsection{An exact criterion for coherent history lifting}
The finite-table criterion can also be sharpened structurally for free
path grammars. Even when each
ordered comparison has a finite native history, can those histories be
chosen compatibly with composition?

Assume $\Gamma$ is a finite acyclic directed multigraph, with vertex set
$\Gamma_0$ and free path category $F\Gamma$. Then $P=\rho(F\Gamma)$ is a finite
partial order on the same vertices. Let
$q:F\Gamma\to P$ be the reachability quotient. A saturated chain is a
chain of cover relations in $P$. We now ask for a section of the
witness-forgetting functor $q$ itself. This is stronger than lifting the
unary identity operation along $\Phi$: that operation already has
$\id_{F\Gamma}$ as a lift. The operation-lifting and witness-selection
questions must not be conflated.

\begin{theorem}[Coherent reachability lifting for free histories]
\label{path-thm:section}
The following are equivalent:
\begin{enumerate}
\item There is an identity-on-objects functor $s:P\to F\Gamma$ with
      $qs=\id_P$.
\item For every comparable pair $a<b$ in $P$, there is exactly one
      saturated chain from $a$ to $b$.
\end{enumerate}
When these hold, choose one graph edge for every cover relation. The
unique saturated chain gives a constructive section by concatenating
those chosen edges. Its histories have length at most $|\Gamma_0|-1$.
\end{theorem}
\begin{proof}
Every cover $a\prec b$ is realized by a graph edge: a path of length
larger than one would give a strict intermediate element. Conversely,
any path from $a$ to $b$ for a cover has length one, since the graph is
acyclic.

Suppose a section exists. On each cover it must choose one such edge.
For any saturated chain $a=a_0\prec a_1\prec\cdots\prec a_k=b$,
functoriality forces the value of the unique order arrow $a\leq b$ to
be the concatenation of those selected cover edges. Two different
saturated chains have different vertex sequences, so these concatenations
are different arrows in a free category. Thus two chains are impossible.

Conversely, choose an edge at each cover. Define $s(a\leq b)$ by the
unique saturated chain, and send identities to empty paths. The
concatenation of the unique saturated chains from $a$ to $b$ and $b$ to
$c$ is a saturated chain from $a$ to $c$, and hence is the unique one.
This proves preservation of composition and the section equation.
A strictly increasing chain visits at most $|\Gamma_0|$ vertices.
\end{proof}

\begin{example}[A finite diamond has individual lifts but no coherent lift]
Take edges $0\to a\to1$ and $0\to b\to1$, with $a\ne b$.
Every arrow of the reachability order has a native path. A section would
have to send $0\leq1$ both to the path through $a$ and to the path through
$b$. They are distinct, so no section exists. This is a failure of
composition-compatible witness recovery, even though all models and all
path choices in this example are finite.
\end{example}

Parallel cover edges do not obstruct existence of a section; one may
choose one. They do obstruct canonical choice under a symmetry swapping
the edges. If a native comparison must preserve every named primitive
edge, the condition is stronger still: $s q$ must fix each generator,
hence every path, and therefore $q$ must already be an isomorphism.
A section that replaces a primitive shortcut by a longer chain is an
ordered realization, not preservation of that primitive's native meaning.

The exact section theorem is a positive class and an obstruction for a
true change from category-valued histories to order-valued existence.
It is not a claim that all Cat-enriched operation families or contextads
are free path categories. For a richer native doctrine, its actual
higher cells, whiskering, endpoint maps, and coherence equations must be
included and checked separately.


\paragraph{Naturality of the positive section.}
For a category of such marked graph presentations, require its maps to
preserve covers and the chosen cover edges. A map then carries the unique
saturated chain to the unique saturated chain between the image endpoints,
so the constructed sections are natural. Without this mark-preservation
condition, a symmetry exchanging parallel cover edges can obstruct a
natural section even though each presentation individually has one. The
condition changes the source contract explicitly; it does not follow from
reachability alone.

\begin{corollary}[Natural operation lifting on marked path presentations]
\label{path-cor:operation-lift}
Let $\mathcal E$ be a finite category of finite acyclic graph presentations
satisfying Theorem~\ref{path-thm:section}, with chosen cover edges and
presentation maps preserving those covers and edges. Give each primitive
graph edge unit cost. Put $C_e=F\Gamma_e$, $P_e=\rho C_e$, and let $q_e:C_e\to P_e$ and
$s_e:P_e\to C_e$ be the natural quotient and section families.
Every natural family of monotone operations $f_e:P_e^n\to P_e$ has a
natural categorical lift
\[
                    \widetilde f_e=s_e f_e q_e^n:C_e^n\to C_e.
\]
It is executable in the finite-table/native-path model of
Proposition~\ref{prop:dc-finite-csp}. Each output arrow is represented by
a path of length at most $|\Gamma_{e,0}|-1$.
\end{corollary}
\begin{proof}
The displayed composite is a functor. Naturality of $q$, $s$ and $f$
gives naturality of $\widetilde f$. Since $q_es_e=\id_{P_e}$,
its reachability image is $f_e$, so it is an actual preimage under the
operation comparison of Theorem~\ref{cat-thm:base}. The graph is finite
and acyclic, so its path category is finite. The resulting functor table
is finite, and the chosen section gives the stated native path for every
output arrow. The path bound is Theorem~\ref{path-thm:section}.
\end{proof}
This construction selects lifts of operations, but does not in general
give a unital section of operation theories: it sends the unary identity
to $s_eq_e$, which need not be $\id_{C_e}$. It also does not recover an
originally executed path. The positive conclusion is precisely the
natural finite-table lift displayed above.

\paragraph{Why groupoid fibre tests do not replace directed lifting.}
For functors between groupoids, equivalence is equivalent to all
iso-comma fibres being nonempty connected groupoids with trivial
automorphism groups. This is a standard full-faithfulness and essential
surjectivity test. Concrete orbit/stabilizer calculations can make it
finite, but it concerns only invertible witnesses. The functor from the
discrete category on $\{0,1\}$ to the walking-arrow category $[1]$ has
terminal iso-comma fibres at both objects and still misses the
noninvertible arrow. It is not an equivalence of categories. Ordinary
categories, reachability preorders, connected components and maximal
groupoids must therefore remain distinct. The source tables and path
composition equations above retain precisely the noninvertible data that
this fibre test misses.

\section{The existing law and possibility example}\label{sec:o7}
\label{subsec:dc-law}

This example is included because it already supplies a positive answer to
part of the same reconstruction-and-transport question in the project's
native contextad lineage. It must not be advertised as a new discovery or
replaced by a weaker generic image criterion. At the monad level, a map
$\pi:L\to D$ gives a reduct of models, and
\eqref{eq:allmodels} identifies the induced operation map at arity $n$
with $\pi_n:Ln\to Dn$. At the raw native level one retains the actual
source-to-observation map and asks which witness functions survive.
These are the two specified levels of the following calculation. The
monad calculation uses all algebra models and their generic arities; it is
not an application of the finite-signature probe theorem. A finite support
test is effective with rational tables or another supplied exact
zero-decision procedure for the coefficients.


Let $\mathsf D X$ be finite probability laws. For effective finite-table
calculations one may uniformly restrict coefficients to nonnegative
rationals; none of the following finite proofs changes. Define
\[
 \mathsf L X=\{(\mu,S):\mu\in\mathsf D X,\quad
     S\subseteq X\text{ finite},\quad\operatorname{supp}\mu\subseteq S\}.
\]
Such $S$ is necessarily nonempty. The chosen set records declared
possible outcomes, not a family of alternative probability laws. For
$k(x)=(\nu_x,S_x)$ the bind is
\begin{equation}\label{eq:dc-Lbind}
 (\mu,S)\mathbin{\gg=}k
 =\left(\sum_x\mu(x)\nu_x,\ \bigcup_{x\in S}S_x\right),
 \qquad \eta(x)=(\delta_x,\{x\}).
\end{equation}
The finite support inclusion, associativity and units follow respectively
from positivity, reassociation of the finite sums, and reassociation of
the unions. This is exactly the basic observation monad in project O7;
O7 also proves its occurrencewise interaction with trees and additive
costs. It is not the convex powerset monad and makes no general
probability--nondeterminism distributive-law claim.

\begin{proposition}[Split forgetting and an exact positive fragment]
\label{prop:dc-law-section}
There are monad morphisms
\[
 \pi:\mathsf L\longrightarrow\mathsf D,\quad(\mu,S)\mapsto\mu,
 \qquad
 s:\mathsf D\longrightarrow\mathsf L,\quad
       \mu\mapsto(\mu,\operatorname{supp}\mu),
 \qquad \pi s=1.
\]
The idempotent $s\pi$ has as its fixed points exactly the pairs with
$S=\operatorname{supp}\mu$. These are exactly the finite outcome data realizable
by finite raw tables all of whose retained rows have strictly positive
weight. Thus a finite support test recognizes the observed image of strictly
positive raw tables. On this fragment $s$ and $\pi$ are inverse at the level of
observed outcome data. This is not an equivalence of raw witness categories,
and general law-and-possibility data are not recovered.
\end{proposition}
\begin{proof}
For a map $f$, positivity gives
$\operatorname{supp}(f_*\mu)=f[\operatorname{supp}\mu]$.
For finite bind it gives
\[
 \operatorname{supp}\left(\sum_x\mu(x)\nu_x\right)
  =\bigcup_{x\in\operatorname{supp}\mu}\operatorname{supp}\nu_x.
\]
These equations and the Dirac law verify naturality and both monad
morphism axioms for $s$; those for $\pi$ are immediate from
\eqref{eq:dc-Lbind}. The fixed-point statement is direct.
A raw table $(I,p,o)$ has $p_i>0$, $\sum_i p_i=1$, and outcome map
$o:I\to X$. Every value of $o$ then has positive aggregate mass, so
$o[I]=\operatorname{supp}(o_*p)$. Conversely the table with
$I=\operatorname{supp}\mu$, $o$ its inclusion, and $p_i=\mu(i)$
realizes any fixed point. This proof also constructs the realizer.
\end{proof}

The same section extends to O7's finite joint cost/tree occurrence
monads: on a joint law $\lambda$ put
$s_K(\lambda)=(\lambda,\operatorname{supp}\lambda)$.
For each addressed tree, the support of the finite independent product
is the product of supports, and support commutes with the graft-and-sum
pushforward. Together with the finite-bind support equation in the proof,
these facts verify the monad-morphism law for $s_K$. Thus the positive
fragment is composition-closed even with dependent tree occurrences and
joint costs. It still forgets zero-weight retained executions and repeated
raw witnesses of the same joint outcome.

The distinction between a \emph{right} inverse of forgetting and recovery
of the original richer semantics is essential. For $X=\{a,b\}$,
\[
 (\delta_a,\{a\})\ne(\delta_a,\{a,b\}),\qquad
 \pi(\delta_a,\{a\})=\pi(\delta_a,\{a,b\}).
\]
Hence no function of the law can recover the original possibility set;
in particular $s\pi\ne1$. The second outcome datum becomes realizable in the
strictly positive target table fragment after forgetting, although it was
not realizable in the strictly positive source table fragment. This is failure of
reflection of the selected native definability, not failure of existence
of a section.

For the larger native repertoire of finite raw tables with nonnegative
weights, charge the number of retained rows. The minimum row cost of
$(\mu,S)$ is exactly $|S|$: its outcome image requires that many rows,
and the table with one row per element of $S$ realizes it. The minimum
row cost of $\mu$ in law semantics is $|\operatorname{supp}\mu|$.
Thus the pair $(\delta_a,\{a,b\})$ needs two retained rows, whereas its
law needs only one. Multiplicities of rows with the same outcome are
still lost by both observations. At the level of reconstructed theories,
\eqref{eq:allmodels} makes these descriptions literal: the operation
objects of arity $n$ are $\mathsf L n$ and $\mathsf D n$, and the
operation comparison is $\pi_n$.

\paragraph{All-path budgets and O7's sharper native-map test.}
Let outcomes carry costs, with $c(a)=\ell$, $c(b)=H$, and
$\ell\leq B<H$. The two displayed pairs have the same joint probability
law, including its cost coordinate, but only the first satisfies
$\forall x\in S,\ c(x)\leq B$. No decision factoring through the law can
answer that all-path question on both objects. Project O7 supplies the
stronger typed version with different causal trees and a real
zero-weight high-cost branch with two continuation holes. The present
finite pair is its smallest information-loss witness, not a replacement
for that typed construction.

O7's presence coefficient semiring is
\[
 E=\{(p,b):p\geq0,\ b\in\{0,1\},\ p>0\Rightarrow b=1\},
 \quad (p,b)+(q,e)=(p+q,b\lor e),
 \quad (p,b)(q,e)=(pq,b\land e).
\]
It is useful to distinguish two maps back from nonnegative reals:
\[
 p\longmapsto(p,[p>0])\quad\text{is a semiring section},
 \qquad p\longmapsto(p,1)\quad\text{is not: it fails at }0.
\]
The first gives $s$ and declares zero-weight rows absent. The second is
the all-present raw-program lift used to retain the original possibilities.
Conflating them would incorrectly claim that either no section exists or
that the section recovers all raw programs.

For an old raw witness map $r:I\to J$ between all-present objects, its
presence balance at a target row is
\[
 1=\bigvee_{i:r(i)=j}1.
\]
It holds exactly when $r$ is surjective. Consequently an old native map
lifts to the presence-enriched native calculus exactly when each input
fibre's witness function is surjective; the mass, body and cost equations
are otherwise unchanged. This is O7's Proposition ``The native
preservation cost''. Its ghost insertion supplies a concrete old map
with no presence-enriched lift. Flag forgetting is faithful on the
retained raw witness maps and generally not full, while subsequent
observation forgets parallel witness functions as well.

Finite-distribution support as a monad morphism is established primary
background, explicitly discussed by Fritz--Perrone--Rezagholi~\cite{DC-FPR}.
The finite proof here does not infer a topological or measurable result,
and does not identify topological support with all declared raw paths.

\section{Finite observations and effective recognition}
\label{mon-sec:limits}
\begin{proposition}[Every finite transcript has a native completion]
\label{mon-prop:finite-transcript}
For any $\alpha\in\Op_U(n)$ and any finite list of queried pairs
$(M_j,\mathbf a_j)$, there is one word $w$ satisfying
$w^{M_j}(\mathbf a_j)=\alpha_{M_j}(\mathbf a_j)$ for every $j$.
\end{proposition}
\begin{proof}
In $\prod_jM_j$ form the submonoid generated by the coordinatewise input
tuples. Its $\alpha$-value is some word in those generators; projecting
and applying naturality yields all the required equalities. For an empty
list any word suffices.
\end{proof}
This also covers any finite list of complete operation tables, since each
finite table has finitely many inputs. It is the elementary form of the
density of $X^*$ in $\widehat{X^*}$. In particular, no finite evaluation
transcript of a nonnative operation can certify nonnativeness among all
natural operations. For $\omega$ the completion can explicitly be chosen
as $x^{k!}$ with $k$ at least the largest queried monoid size.

A deterministic black-box procedure that claims nonnativeness after
finitely many evaluation queries would make the same queries and receive
the same answers for that completing word, including when its queries
are adaptive. It therefore cannot be sound. This does not contradict the
positive certificates: a finite-word point can be isolated, whereas a
nonword point cannot be isolated from all finite words.

\begin{theorem}[Undecidability under a strong effectiveness promise]
\label{mon-thm:undecidable}
There is no algorithm that, given code for a promised total computable
natural unary operation on all finite monoids, always decides whether it
has a finite native monoid-term representative. There is also no sound
procedure on such codes that halts on exactly the nonnative operations.
\end{theorem}
\begin{proof}
Encode a finite monoid by a positive integer $m$, carrier
$\{0,\ldots,m-1\}$, identity $0$, and an explicit $m$-by-$m$
multiplication table. The monoid axioms and well-formedness of this
encoding are decidable by finite inspection; malformed inputs can be
rejected. An input element is its carrier index. Given a machine $e$, use
a step convention in which any halting time is an integer $t\geq1$.
Build this evaluation program: on a finite monoid
$M$ of size $m$ and $a\in M$, simulate $e$ for $m$ steps. If it halts
at time $t\leq m$, output $a^{t!}$; otherwise output $a^\omega$ using
the finite-table algorithm. This program is total uniformly in $e$.

If $e$ never halts, the resulting family $\alpha^e$ is exactly $\omega$.
If $e$ halts at time $t$, it is exactly the word operation $x^{t!}$:
for $m\geq t$ this is the specified branch, and for $m<t$ Lemma
\ref{mon-lem:factorial} shows that $a^{t!}=a^\omega$. Thus in either case
the family is natural, not merely a partially verified candidate. It is
native exactly when $e$ halts, by Corollary~\ref{mon-cor:length}.
A nativeness decider would decide halting; a semidecision for nonnativeness
would semidecide nonhalting. Both are impossible.
\end{proof}
The theorem concerns a uniform input-code problem, not the difficulty of
particular named operations. It does not preclude a compiler that succeeds
on a specified decidable sublanguage, one supplied with a valid length
bound, or the positive semidecision already given. The halting reduction
is recorded as an elementary consequence, without a priority claim.


\subsection{What an implementation claim additionally needs}
There are several distinct access contracts.
\begin{enumerate}
\item \emph{Ideal semantic access} supplies a natural operation on an
entire admitted model category. Naturality includes the product projections,
generated-subalgebra inclusions and reduct maps used in the proofs.
\item \emph{Exact finite-probe access} supplies the output element on a
specified finite algebra. The theorem's single probe may have an enormous
table. A physical interface exposing only selected measurements does not
provide that element without a justified separating readout family.
\item \emph{Finite observational access} supplies only a transcript. It
neither proves global naturality nor licenses the unobserved test models.
The finite-interpolation theorem explains its specific limitation here.
\item \emph{Source execution access} supplies a verified implementation of
each primitive, its typing and resource contract, and any required reset,
branching, equality test, iteration or state inspection. It is independent
of the existence of an ambient semantic morphism.
\end{enumerate}
The primitive interpretation in Theorem~\ref{alg-thm:compilerimage} is
supplied, but its source term is \emph{extracted}; no global decompiler is
assumed. The finite-table lift in Proposition~\ref{prop:dc-finite-csp}
returns an executable table in its declared machine model, not a term in
an unrelated language. The finite-probe search terminates under a native
bound and semidecides a positive result without one; it is not a general
negative decision procedure. These distinctions are part of the theorem
statements, not qualifications that can be omitted when applying them.

\section{What this establishes and what remains open}

\subsection{A common reconstruction architecture}
The mathematical output is one architecture with distinct obligations.
Given native grammar and semantic presentations, construct $\Op_d$ and
the interpretation $e_d$. Given a doctrine change, construct the actual
reduct and arity comparison and hence $\Phi_u$. Prove native preservation
from primitive translation. Then use an independent recognition mechanism,
rather than the definition of the image, to test reflection or lifting.

In the finite algebraic class, singleton-fibre probes and finite source-test
coverage provide such a mechanism. The compiler-image test is a finite
generated-subalgebra calculation, and the conservative transport theorem
extracts a source term for an already related semantic operation. Its
coverage condition and comparison cohere under composition. Finite monoids
instantiate every part and identify the exact completion. The
semigroup-to-monoid reduct is a positive doctrine change; restriction to
groups shows how a hypothesis can fail.

For the genuinely enriched change $\Cat\to\Preord$, the same operation
comparison can lose arrow actions and fail to cover even a binary target
operation on finite presentations. The finite source-table constraints
give a complete recognition algorithm in their stated machine model; the
acyclic free-path criterion supplies a structural positive class for
composition-compatible lifting. Neither construction silently reinstates
the discarded witness identity. The law-and-possibility example uses a
monad map to realize the same comparison pattern and shows how a
composition-preserving section can recover an exact fragment without
recovering the richer original executions.

\subsection{Credit to earlier results at their actual strength}
Lucyshyn-Wright--Parker Proposition~4.34 and Theorem~4.35 are the closest
general positive recognition results. Bourke--Garner provide a related
pretheory--monad adjunction and presentation/density criteria
\cite{BourkeGarner}. Reiterman-style profinite theory identifies the
finite-model completion; isolated finite words and the finite-carrier
$\operatorname{Pol}$--$\operatorname{Inv}$ mechanism already supply
recognition in their domains \cite{ACMU,Costa,Geiger}.
Baez--Williams organize changes of enrichment, theory and models
\cite{BW}. Finite-power/subalgebra/quotient realizations also have an
established topological-clone theory; its continuity hypotheses and
resource scope should not be dropped \cite{BodirskyPinsker}.
These are alternatives and precedents to apply where their hypotheses
hold, not obstacles manufactured by the present framework.

The earlier CALF and finite linear formation results already reconstruct
native operational quotients and provide bounded algorithms. The project's
O1 criterion derives mixed structure from independently specified cartesian,
occurrence and Beck--Chevalley data; O5 lifts local full-state certificates
through its admitted finite source language. O7 already proves the
law-and-possibility monad, a native paired-history span criterion, the exact
surjectivity test for lifting all-present witness maps, and coherent local
observation reconstruction \cite{CALF,Formation,O1,O5,O7}.
In O7 the comparisons
\[
                  R(g)\circ R(f)\longrightarrow R(g\circ f)
\]
can aggregate histories and are generally noninvertible. A canonical
ambient representative is not thereby generated by the original source
language. The support-minimal section discussed here answers a different
question from O7's all-present lift and its original observation reflection.

No theorem in this note identifies an arbitrary Cat-enriched operation
theory with the project's native relative modules. Such modules can have
different endpoint bases and require their own units, multiplications,
pure maps, endpoint transport and coherence cells. A strong comparison
requires the appropriate invertible coherences; a lax comparison must state
its compositor direction. A groupoid calculation on an independently
specified witness doctrine is not automatically an O1 or O7 instance.

\subsection{Answer to the original ambition}
This goes beyond a fixed finitary typed linear framework in two proved
ways: it treats an entire category of finite algebraic presentations with
an explicit semantic completion and source-term extraction across grammar
changes, and it analyzes a genuinely witness-erasing enriched doctrine
change with positive and negative lifting results. It is not just a list
of unrelated models, because each case specifies the same native map,
reconstructed operations and doctrine comparison and tests one of their
properties.

It does not solve arbitrary reconstruction of a native repertoire from
unspecified observations. Arbitrary equations may destroy finite
singleton recognizers. Different arities need not match. A target semantic
operation need not be in the comparison image. Individual lifts need not
be natural or composition-compatible. A natural ambient realization need
not be a generated source program. Finite evidence does not by itself
validate global naturality or physical availability. The counterexamples
show that these are mathematical obstructions, not missing terminology.

The remaining research problem is therefore specific: for another
independently meaningful native doctrine, construct a checkable generic
probe or invariant-based recognizer from its primitive data, identify the
actual reduct and arities, and prove source-program and witness transport
with the intended resource and access contract. A theorem doing this for
a new native class, or deriving a contextad instance without assuming its
realization property, would be a substantive further result. The present
note supplies a proved conditional framework and several exact classes;
it makes no claim that this further universal result has been obtained.

\begin{thebibliography}{99}
\bibitem{LWP}
R. B. B. Lucyshyn-Wright and J. Parker.
\emph{Enriched structure--semantics adjunctions and monad--theory
equivalences for subcategories of arities}.
Theory and Applications of Categories 41 (2024), 1873--1918.
Proposition 4.24, Theorem 4.27, Proposition 4.34 and Theorem 4.35.
\url{https://arxiv.org/html/2305.07076}.
\bibitem{BourkeGarner}
J. Bourke and R. Garner. \emph{Monads and theories}.
Advances in Mathematics 351 (2019), 1024--1071.
Theorems 6, 38 and 43.
\url{https://arxiv.org/html/1805.04346}.
\bibitem{ACMU}
J. Ad\'amek, L.-T. Chen, S. Milius and H. Urbat.
\emph{Reiterman's Theorem on Finite Algebras for a Monad}.
Author-hosted January 2021 manuscript, Definition 4.5, Construction 4.6,
Remark 4.7, Definition 5.16, Example 5.17 and Lemma 5.19.
\url{https://www8.cs.fau.de/wp-content/uploads/media/research/papers/reitermanstheoremonfinitealgebrasforamonad.pdf}.
\bibitem{Costa}
A. Costa. \emph{A profinite approach to complete bifix decodings of
recurrent languages}. Coimbra Preprint 22--25 (2022), Section 2.3.
\url{https://www.mat.uc.pt/preprints/ps/p2225.pdf}.
\bibitem{Geiger}
D. Geiger. \emph{Closed systems of functions and predicates}.
Pacific Journal of Mathematics 27 (1968), 95--100.
\url{https://msp.org/pjm/1968/27-1/pjm-v27-n1-p08-s.pdf}.
\bibitem{BodirskyPinsker}
M. Bodirsky and M. Pinsker. \emph{Topological Birkhoff}.
Theorems 4 and 14; primary author manuscript.
\url{https://www.dmg.tuwien.ac.at/pinsker/papers/topobirkhoff/pinsker_topobirkhoff.pdf}.
\bibitem{BW}
J. C. Baez and C. Williams.
\emph{Enriched Lawvere Theories for Operational Semantics}.
EPTCS 323 (2020), 106--135. Theorem 8 and Sections 5--6.
\url{https://arxiv.org/html/1905.05636v3}.
\bibitem{DC-FPR}
T. Fritz, P. Perrone and S. Rezagholi.
\emph{Probability, valuations, hyperspace: Three monads on Top and the
support as a morphism}. Mathematical Structures in Computer Science 31
(2021), 850--897.
\url{https://doi.org/10.1017/S0960129521000414}.
\bibitem{CALF}
\emph{When observations warrant new structure}.
Earlier programme research note, 3 November 2026.
\bibitem{Formation}
\emph{Operational Theory Formation Through Interaction}.
Programme research note, 5 November 2026.
\bibitem{O1}
\emph{General Independent Interaction Criterion}.
Project O1 dossier. Main theorem and marked native recognition theorem.
\bibitem{O5}
\emph{Contextual Refinement and Controlled Recovery}.
Project O5 dossier. Source-language compositional lifting theorem.
\bibitem{O7}
\emph{Law and Possibility Repair for Dependent Protocols}.
Project O7 dossier, 1 November 2026. Main theorem, native preservation
criterion and local observation reflection.
\end{thebibliography}
\end{document}
